For \(A\in\mathcal C_n\), abbreviate \(q_{\sqrt{s}}(A;S)\) by \(q_s(A)\) and \(v^*(\sqrt{s};S)\) by \(v(s)\).  If \(u\in H_n\) is a unit vector, put \(t(u)=u^\top P_Su\).  We first eliminate the continuum of quality-indexed spherical constraints by proving
\[
\sup_{0\leq s\leq1}\{q_s(A)-v(s)\}
=\sup_{\substack{u\in H_n\\\|u\|_2=1}}
\{u^\top Au-v(t(u))\}.
\tag{1}
\]
For the lower bound in \((1)\), fix \(u\) and take \(s=t(u)\); then \(q_s(A)\geq u^\top Au\).  For the upper bound, let \(u_s\) attain \(q_s(A)\) and put \(t=t(u_s)\geq s\).  The oracle value is nonincreasing because its adversary class shrinks as the quality lower bound increases, so \(v(s)\geq v(t)\).  Therefore
\[
q_s(A)-v(s)=u_s^\top Au_s-v(s)
\leq u_s^\top Au_s-v(t),
\]
which is bounded by the right side of \((1)\).

Now enumerate sign-pair representatives and write every \(B\in\mathcal C_n\) as \(A(b)\), \(b\in\Delta_m\).  The cone-lift scalar dual and the oracle minimization give, for \(t\in[0,1]\),
\[
v(t)=\inf_{\substack{b\in\Delta_m,\ \eta\geq0,\ \ell\in\mathbb R\\
\ell I_{n-1}-Q_H^\top\{A(b)+\eta P_S\}Q_H\succeq0}}
\{\ell-\eta t\}.
\tag{2}
\]
The matrix inequality says exactly that \(\ell\) dominates the largest eigenvalue on \(H_n\).  At \(t=1\), the infimum in \((2)\) may occur only as \(\eta\to\infty\); retaining an infimum is therefore essential.  Substituting \((2)\) into \((1)\) changes its negative infimum into a supremum and yields exactly the robust semidefinite representation in the theorem.  Taking the epigraph gives the displayed finite constraints.

Here is an explicit finite breakpoint procedure.  From a centered full-rank score basis, compute \(P_S\) and an orthonormal \(Q_H\), and enumerate the finite balanced sign-pairs.  Replace the sole matrix inequality in \((2)\) by nonnegativity of all principal minors.  Together with \(\|u\|_2^2=1\), the linear simplex equalities, and coordinate nonnegativity, this is a finite list of polynomial equalities and inequalities whose coefficients are entries of the score projector and the enumerated signs.  Recursively partition each coordinate line at the real roots of these polynomials and of their successive discriminants and resultants.  On each resulting sign-invariant cell, record the active simplex face and the eigenvalue branch determined by
\[
\det\!\left(\ell I_{n-1}-Q_H^\top\{A(b)+\eta P_S\}Q_H\right)=0.
\]
There are finitely many input polynomials of finite degree, so this recursive decomposition has finitely many cells after identically zero polynomials are passed to the next lower-dimensional stratum.  The projected cell endpoints in \(t=u^\top P_Su\), together with \(t=0,1\), are the required breakpoints.  Unbounded \(\eta\) is compactified by \(h=(1+\eta)^{-1}\); the cell \(h=0\) contributes only when \(t=1\), and its finite objective limit is the compressed eigenvalue \(\lambda_{\max}(P_SA(b)P_S\mid S)\), exactly as in the boundary part of the cone-lift proof.  Optimizing the polynomial objective over every cell and comparing finitely many cell optima computes the robust supremum and then the outer minimum.  For numerical score input, interval root isolation and semidefinite bounds give the stated \(\varepsilon\)-solution.  This construction uses only the score basis, \(n\), and the balanced signs.

It remains to establish attainment and the converse.  The function \(v(t)\) is continuous on \([0,1]\).  Indeed, for a fixed covariance, rotate the \(S\)- and \(S^\perp\)-components of a feasible unit vector to change its score energy; the Euclidean displacement tends uniformly to zero with the change in score energy, and
\[
|u^\top Au-w^\top Aw|\leq2\|A\|_{\mathrm{op}}\|u-w\|_2\leq2n\|u-w\|_2,
\]
because \(A\succeq0\) and \(\operatorname{tr}(A)=n\).  This proves uniform continuity of each \(q_s(A)\), uniformly over \(A\in\mathcal C_n\), including the endpoints; minimization over the compact covariance hull preserves it.  Equation \((1)\) then expresses \(\Delta_{\mathrm{reg}}(A;S)\) as the maximum of a jointly continuous function over the compact unit sphere.  It is continuous in \(A\).  Since \(\mathcal C_n\) is compact, an outer minimizer \(A_{\mathrm{reg}}\) exists.  Balanced-covariance realizability converts a convex representation of this matrix into a symmetrized \(d_{\mathrm{reg}}\in\mathcal D_n\).

The exact risk theorem gives, for every \(d\in\mathcal D_n\),
\[
\sup_{\rho\in[0,1]}
\left\{
\sup_{\mu\in\mathcal Y(\rho,M;S)}
\mathbb E_d[(\widehat\tau_d-\tau_n)^2]-V^*(\rho;S)
\right\}
=\frac{4M}{n}\Delta_{\mathrm{reg}}(A_d;S)
\geq R^*(n,S).
\]
Equality holds at \(A_{\mathrm{reg}}\), proving both attainability and the matching lower bound against every quality-agnostic design.

Finally, at \(\rho=0\), \(q_0(A;S)=\lambda_{\max}(A\mid H_n)\).  The trace bound gives \(q_0(A;S)\geq c_n\), while complete randomization has covariance \(c_nP_H\) and achieves equality.  It is therefore a \(\rho=0\) oracle design.  Its support is all of \(\mathcal Z_n\).  Any symmetrized representation of \(A_{\mathrm{reg}}\) uses only assignments in \(\mathcal Z_n\), so its support is contained in the support of this single oracle design.  This answers the final support question affirmatively, with one oracle support already sufficient.